\section{Conclusions}\label{sec:conclusions}

In future work, we plan to apply the Bayesian tag-and-probe method to
tau lepton decays. The main limitation of this approach is its reliance
on the Sudakov form factors of the hard function of the likelihood
ratios of the BRST-invariant differential cross section. Remarkably, the
massive N-jettiness slicing scheme encodes the rapidity gap fraction for
any number of emissions. We leave the study of the N-subjettiness ratios
of the holistic azimuthal decorrelation at two loops to future work. The
implications for physics beyond the Standard Model are profound.

